% from ega0-16-fr.tex
\oldpage[0\textsubscript{IV-1}]{26}
Soit $(M_n)$ une filtration $\mathfrak{q}$-bonne de $M$ (Bourbaki,
\emph{Alg.~comm.}, chap.~III, \S~3, no~1)~; alors
$\operatorname{long}(M/M_n)$ est égal, pour $n$ grand, à un polynôme en $n$
dont le degré et le coefficient dominant sont les mêmes que ceux de
$P_{\mathfrak{q}}(M,n)$.
\end{lemma}

\begin{proof}
En effet, il existe $n_0$ tel que $M_{n+1}=\mathfrak{q}M_n$ pour $n\geq n_0$,
d'où les inclusions
$\mathfrak{q}^{n+n_0}M\subset M_{n+n_0}=\mathfrak{q}^nM_{n_0}
\subset\mathfrak{q}^nM\subset M_n$ pour $n$ grand, ce qui entraîne
\[
  P_{\mathfrak{q}}(M,n+n_0)\geq
  \operatorname{long}(M/M_{n+n_0})\geq
  P_{\mathfrak{q}}(M,n)\geq\operatorname{long}(M/M_n)
\]
et par suite la conclusion.
\end{proof}

Si $\mathfrak{q}'$ est un second idéal de définition, on a
$\mathfrak{q}'\supset\mathfrak{q}^m$ pour un entier $m$, donc
$P_{\mathfrak{q}'}(M,n)\leq P_{\mathfrak{q}}(M,mn)$~; par suite le degré de
$P_{\mathfrak{q}}(M,n)$ ne dépend pas de l'idéal de définition
$\mathfrak{q}$ considéré~; on le note $d(M)$.

\begin{lemma}[16.2.2.2]
\label{0.16.2.2.2-fr}
Soit $0\to M'\to M\to M''\to0$ une suite exacte de $A$-modules de type
fini. Alors le polynôme
\[
  P_{\mathfrak{q}}(M,n)-P_{\mathfrak{q}}(M',n)
  -P_{\mathfrak{q}}(M'',n)
\]
a un degré $\leq d(M')-1\leq d(M)-1$.
\end{lemma}

\begin{proof}
En effet, la filtration des
$M'_n=M'\cap\mathfrak{q}^nM'$ est $\mathfrak{q}$-bonne (Bourbaki,
\emph{Alg.~comm.}, chap.~III, \S~3, no~1, prop.~1)~; comme
\[
  \operatorname{long}(M/\mathfrak{q}^nM)
  =\operatorname{long}(M''/\mathfrak{q}^nM'')
  +\operatorname{long}(M'/M'_n),
\]
le lemme résulte aussitôt de (16.2.2.1) appliqué à $M'$.
\end{proof}

\begin{theorem}[16.2.3]
\label{0.16.2.3-fr}
\textup{(Krull-Chevalley-Samuel).}
Soient $A$ un anneau semi-local noethérien, $\mathfrak{r}$ son radical,
$M\neq0$ un $A$-module de type fini, $d(M)$ le degré du polynôme de
Hilbert-Samuel de $M$ (pour un idéal de définition quelconque de $A$)~; soit
d'autre part $s(M)$ la borne inférieure des entiers $n$ tels qu'il existe des
éléments $x_1,\ldots,x_n$ de $\mathfrak{r}$ pour lesquels
$M/(x_1M+\cdots+x_nM)$ soit de longueur finie. Alors on a
\[
  d(M)=s(M)=\dim(M).
\]
\end{theorem}

\begin{lemma}[16.2.3.1]
\label{0.16.2.3.1-fr}
Pour tout $x\in\mathfrak{r}$, soit ${}_xM$ le sous-module de $M$ formé des
éléments annulés par $x$. Alors~:
\begin{enumerate}
  \item[\rm{(i)}] On a $s(M)\leq s(M/xM)+1$.
  \item[\rm{(ii)}] Soient $\mathfrak{p}_i$ les idéaux premiers appartenant à
  $\operatorname{Supp}(M)$ et tels que
  \[
    \dim(A/\mathfrak{p}_i)=\dim(M)\quad(1\leq i\leq m).
  \]
\end{enumerate}
Si $x\notin\mathfrak{p}_i$ pour tout $i$, on a
$\dim(M/xM)\leq\dim(M)-1$.
\begin{enumerate}
  \item[\rm{(iii)}] Si $\mathfrak{q}$ est un idéal de définition de $A$, le
  polynôme
  \[
    P_{\mathfrak{q}}({}_xM,n)-P_{\mathfrak{q}}(M/xM,n)
  \]
  est de degré $\leq d(M)-1$.
\end{enumerate}
\end{lemma}

\begin{proof}
L'assertion (i) est triviale, car si $(x_i)_{1\leq i\leq m}$ est tel que pour
le module $N=M/xM$, $N/(x_1N+\cdots+x_mN)$ soit de longueur finie, il suffit
d'observer que ce module est isomorphe à
$M/(xM+x_1M+\cdots+x_mM)$.

Pour prouver (ii), on remarque que, si $\mathfrak{a}$ est l'annulateur de
$M$, les idéaux premiers qui contiennent l'annulateur de $M/xM$ sont ceux qui
contiennent $\mathfrak{a}+Ax$ ($0_{\mathrm I}$, 1.7.5)~; aucun des
$\mathfrak{p}_i$ ne peut donc contenir $\mathfrak{a}+Ax$ et par définition de
$\dim(M)$ (16.1.7) toute chaîne d'idéaux premiers contenant
$\mathfrak{a}+Ax$ a donc une longueur $\leq\dim(M)-1$.

